% !TEX program = lualatex
% !TEX root = poster.tex
% Unofficial University of Cambridge Poster Template
% https://github.com/andiac/gemini-cam
% a fork of https://github.com/anishathalye/gemini
% also refer to https://github.com/k4rtik/uchicago-poster

\documentclass[final]{beamer}

% ====================
% Packages
% ====================

\usepackage[T1]{fontenc}
\usepackage{lmodern}
\usepackage[size=custom,width=91.44,height=60.96,scale=1.0]{beamerposter}
\usetheme{gemini}
\usecolortheme{upenn}
\usepackage{graphicx}
\usepackage{booktabs}
\usepackage[numbers]{natbib}
\usepackage{tikz}
\usepackage{pgfplots}
\pgfplotsset{compat=1.14}
\usepackage{anyfontsize}
\usepackage{amsmath}
\usepackage{amssymb}

% Tighten caption spacing
\setlength{\abovecaptionskip}{0.3em}
\setlength{\belowcaptionskip}{0.1em}

% ====================
% Lengths
% ====================

% If you have N columns, choose \sepwidth and \colwidth such that
% (N+1)*\sepwidth + N*\colwidth = \paperwidth
\newlength{\sepwidth}
\newlength{\colwidth}
\setlength{\sepwidth}{0.02\paperwidth}
\setlength{\colwidth}{0.31\paperwidth}

\newcommand{\separatorcolumn}{\begin{column}{\sepwidth}\end{column}}

% ====================
% Title
% ====================

\title{Conditional Laws for the Maximal Values of Anchored Random Walk Functionals}

\author{Jonathan Pipping \and Jiaoyang Huang \and Abraham J. Wyner}

\institute[shortinst]{University of Pennsylvania, Department of Statistics \& Data Science}

% ====================
% Footer (optional)
% ====================

\footercontent{
  \href{https://www.upenn.edu}{University of Pennsylvania} \hfill
  2025 PhD Poster Presentation \hfill
  \href{mailto:jpipping@wharton.upenn.edu}{jpipping@wharton.upenn.edu}}

% ====================
% Body
% ====================

\begin{document}

\begin{frame}[t]
\begin{columns}[t]
\separatorcolumn

\begin{column}{\colwidth}

  \begin{block}{Problem Setup}

    \textbf{Setup:} Let $(X_k)_{k=0}^N$ be a random walk with $X_0=0$ and i.i.d. increments $\Delta X_k$ of mean $\mu$ and variance $\sigma^2$, adapted to $(\mathcal F_k)$. For a position $d$ at time $k$, define the functional
    \begin{equation*}
        f(k, X_k) = \mathbb{P}(X_N > 0 \mid X_k = d)
    \end{equation*}
    Then the maximum of $f$ over the time interval $[0,N]$ is
    \begin{equation*}
        M = \max_{0\le k\le N} f(k,X_k).
    \end{equation*}
    \textbf{Question:} What is the law of $M$ conditional on the event $X_N<0$? A closed-form solution is difficult to obtain, so we consider the limiting case.

    \textbf{Brownian Limit:} By Donsker's invariance principle, as $N\to\infty$,
    \begin{equation*}
        \frac{X_{\lfloor Nt \rfloor}}{\sigma\sqrt{N}} \xrightarrow{d} B_t + \mu^* t, \qquad \mu^* = \frac{\sqrt{N}\mu}{\sigma},
    \end{equation*}
    In the limit,
    \begin{equation*}
        f(t,B_t) = \mathbb{P}(B_1 \ge -\mu^* \mid \mathcal F_t) = \Phi\!\left(\frac{B_t+\mu^*(1-t)}{\sqrt{1-t}}\right),
    \end{equation*}
    This represents a bounded martingale with $f(0,B_0)=\Phi(\mu^*)$ whose terminal value is in $\{0,1\}$. The analogous maximum over the time interval $[0,1]$ is
    \begin{equation*}
        M = \sup_{0\le t < 1} f(t,B_t),
    \end{equation*}
    and we study the law of $M$ conditional on the event $B_1 < -\mu^*$.
  \end{block}

  \begin{block}{Symmetric Random Walks}

    \textbf{Functional:} Since $\mu = 0$, the functional simplifies to
    \begin{equation*}
        f(t,B_t) = \mathbb{P}(B_1 \ge 0 \mid \mathcal F_t) 
        = \Phi\!\left(\frac{B_t}{\sqrt{1-t}}\right),
    \end{equation*}
    \textbf{Unconditional Distribution:} Optional stopping at the first hitting time of level $x$ gives the unconditional distribution of $M$:
    \begin{equation*}
        F_M(x) = 
        \begin{cases}
            0, & x < \tfrac12 \\[6pt]
            1 - \tfrac{1}{2x}, & \tfrac12 \le x < 1 \\[6pt]
            1, & x = 1
        \end{cases}
    \end{equation*}
    \textbf{Conditional Distribution:} The conditional distribution of $M$ given $B_1<0$ is
    \begin{equation*}
        F_{M \mid B_1 < 0}(x) = 2 - \frac{1}{x}, \quad x \in [\tfrac12, 1]
    \end{equation*}
    \textbf{Density Function}: Its corresponding density function is
    \begin{equation*}
        f_{M \mid B_1 < 0}(x) = \frac{1}{x^2}, \quad x \in (\tfrac12, 1)
    \end{equation*}

  \end{block}

\end{column}

\separatorcolumn

\begin{column}{\colwidth}

  \begin{block}{Asymmetric Random Walks}

    \textbf{Functional:} In the presence of drift ($\mu^* \neq 0$), the functional
    \begin{equation*}
        f(t,B_t) = \mathbb{P}(B_1 \ge -\mu^* \mid \mathcal F_t)
        = \Phi\!\left(\frac{B_t+\mu^*(1-t)}{\sqrt{1-t}}\right),
    \end{equation*}
    is a bounded martingale with $f(0,B_0)=\Phi(\mu^*)$.

    \textbf{Unconditional Distribution:} Optional stopping at the first hitting time of level $x$ gives the distribution of $M$:
    \begin{equation*}
        F_M(x) = 
        \begin{cases}
            0, & x < \Phi(\mu^*) \\[6pt]
            1 - \tfrac{\Phi(\mu^*)}{x}, & \Phi(\mu^*) \le x < 1 \\[6pt]
            1, & x = 1
        \end{cases}
    \end{equation*}
    \textbf{Conditional Distribution:} Conditional on $B_1< -\mu^*$,
    \begin{equation*}
      F_{M \mid B_1 < -\mu^*}(x) = 1 - \dfrac{\Phi(\mu^*)}{\Phi(-\mu^*)}\cdot\dfrac{1-x}{x}, \quad x \in [\Phi(\mu^*),1]
    \end{equation*}
    \textbf{Density Function}:
    \begin{equation*}
      f_{M \mid B_1 < -\mu^*}(x) = \frac{\Phi(\mu^*)}{\Phi(-\mu^*)}\cdot\frac{1}{x^2},
      \quad x \in (\Phi(\mu^*),1).
    \end{equation*}
  \end{block}

  \begin{block}{Complementary Processes}

    \textbf{Definition:} Consider the complementary martingales
    \begin{equation*}
      p_t = \mathbb{P}(B_1 \ge -\mu^* \mid \mathcal F_t),
      \qquad q_t = 1-p_t = \mathbb{P}(B_1 < -\mu^* \mid \mathcal F_t).
    \end{equation*}

    \textbf{Terminal Behavior:} At time $1$, exactly one of $\{p_t,q_t\}$ converges to $1$, while the other converges to $0$ (a.s.). Our object of interest is defined as
    \begin{equation*}
      M = \sup_{0 \le t < 1} \big( p_t \,\mathbf 1_{\{p_1=0\}} \;+\; q_t \,\mathbf 1_{\{q_1=0\}} \big),
    \end{equation*}
    i.e.\ the running supremum of whichever process ultimately vanishes.

    \textbf{Branch Distributions:} For a process that starts at $s \in (0,1)$ and is conditioned to terminate at $0$, denote by $F_s(x)$ the distribution function of its maximum:
    \begin{equation*}
      F_s(x) = \mathbb{P}(\sup_{0\le t<1} p_t \le x \mid p_0=s,\,p_1=0)
    \end{equation*}

    \textbf{Distribution of the Eventual Vanishing Process:} Let $p_0 := \Phi(\mu^*)$. Then the law of the eventual vanishing process is the mixture
    \begin{equation*}
      F_{\mathrm{vanish}}(x)
      \;=\;
      \Phi(\mu^*)\, F_{\,1-\Phi(\mu^*)}(x)
      \;+\;
      (1 - \Phi(\mu^*))\,F_{\,\Phi(\mu^*)}(x)
    \end{equation*}
    with support $x \in [\min\{\Phi(\mu^*),1-\Phi(\mu^*)\},\,1]$.
\end{block}


 

\end{column}

\separatorcolumn

\begin{column}{\colwidth}

  \begin{block}{Monte Carlo Simulation}

    \textbf{Validation:} We validate our theoretical results with Monte Carlo simulations. The symmetric case is shown in Figure \ref{fig:symmetric}.
    \begin{figure}
      \centering
      \includegraphics[width=0.9\textwidth]{figures/symmetric.png}
      \caption{Symmetric case: simulated distribution (red) vs. theoretical density (blue).}
      \label{fig:symmetric}
    \end{figure}
    \vspace{-0.75em}
  \end{block}

  \begin{block}{Application to In-Game Win Probabilities}

    \textbf{NFL Win Probabilities:} Win probability trajectories from closely-matched games approximately follow our theoretical symmetric distribution.
    
    \begin{figure}
      \centering
      \includegraphics[width=0.9\textwidth]{figures/nfl.png}
      \caption{Distribution of maximal win probabilities of eventual NFL losers (2002-2024).}
    \end{figure}
    \vspace{-0.75em}
  \end{block}

  \begin{block}{References}

    \nocite{*}
    \footnotesize{\bibliographystyle{apalike}\bibliography{poster}}

  \end{block}

\end{column}

\separatorcolumn
\end{columns}
\end{frame}

\end{document}
